\section{强化学习范式}

\gls{rl}的主要目标是构想一个\emph{智能体（agent）}，通过与其\emph{环境（environment）}交互学习最大化某种效用度量。交互过程建模为序贯决策问题。每一步，智能体在环境中执行一个\emph{动作（action）}。作为该动作和外部因素的结果，环境转换到新\emph{状态（state）}并向智能体提供强化信号。该信号称为\emph{奖励（reward）}，对动作质量提供反馈。智能体的目标是最大化累积奖励和——即\emph{回报（return）}。环境遵循的序贯过程（包括转移动态和奖励）由\glsfirst{mdp}建模。为最大化奖励，智能体必须自主平衡探索（exploration）和利用（exploitation）。它需探索\gls{mdp}的新区域以防止错过机会，同时考虑利用已知有利区域以最大化效用。

\gls{rl}框架受心理学关于人类学习过程发现的启发。它足够通用以涵盖许多问题，特别是\gls{ml}的其他子领域——\emph{监督学习（supervised learning）}和\emph{无监督学习（unsupervised learning）}。它也足够具体以产生有意义的理论结果指导算法设计。

在实证方面，框架的有效性由其带来的突破性进展所支持。其中部分成就包括：
机器人运动\cite{haarnoja2018learning}；
自动驾驶\cite{kiran2021deep}；
在专家级别玩复杂视频游戏\cite{vinyals2019grandmaster}；
设计性能更高且能耗更低的计算机芯片\cite{mirhoseini2021graph}；
控制托卡马克中等离子体形状\cite{degrave2022magnetic}。

然而，大多数这些应用需要对原始\gls{rl}框架进行轻微扩展，例如考虑状态的部分可观测性。原始框架的另一个限制可能是动力系统中存在延迟。下一节将解释为何考虑延迟是实现出色实证结果的关键。尽管\gls{ml}和\gls{rl}研究进展如此迅速，但我们建议读者花些时间深入探讨延迟问题。
